% 自动生成：请勿手动编辑。来源：图片/_gen_exp_figs.py
% 场景07 SL — 子图 b：MIKU
\def\datapath{data/07_narrow_cones_cmp/}
\pgfplotsset{
    discard if not/.style 2 args={x filter/.code={
        \edef\tempa{\thisrow{#1}}\edef\tempb{#2}
        \ifx\tempa\tempb\else\def\pgfmathresult{nan}\fi
    }},
    discard if not2/.style n args={4}{x filter/.code={
        \edef\tempa{\thisrow{#1}}\edef\tempb{#2}
        \edef\tempc{\thisrow{#3}}\edef\tempd{#4}
        \ifx\tempa\tempb
            \ifx\tempc\tempd\else\def\pgfmathresult{nan}\fi
        \else\def\pgfmathresult{nan}\fi
    }}
}
\begin{tikzpicture}
\begin{axis}[
    width=0.9\linewidth, height=4.2cm,
    xlabel={$s$ (m)}, ylabel={$l$ (m)},
    xmin=0, xmax=50.0, ymin=-2.2, ymax=2.2,
    grid=both, grid style={gray!25},
    axis line style={thick},
    tick label style={font=\scriptsize},
    label style={font=\footnotesize},
    legend style={at={(0.98,0.02)}, anchor=south east, font=\scriptsize, draw=none, fill=white, fill opacity=0.85},
]
\addplot[fill=bglight, draw=none, forget plot, on layer=axis background] coordinates {(0,-1.875) (50.0,-1.875) (50.0,1.875) (0,1.875)} \closedcycle;

\draw[black!55, semithick] (axis cs:0,-1.875) -- (axis cs:50.0,-1.875);
\draw[black!55, semithick] (axis cs:0,1.875) -- (axis cs:50.0,1.875);
\draw[black!30, dashed, thin] (axis cs:0,0) -- (axis cs:50.0,0);
% 可行带阴影（blocked 时按 blocked_s 截断）
\addplot[name path=lmin_miku, draw=vibrantorange, thin, dashed, forget plot, ]
    table[col sep=comma, x=s, y=l_min, discard if not={mode}{miku}] {\datapath sl.csv};
\addplot[name path=lmax_miku, draw=vibrantorange, thin, dashed, forget plot, ]
    table[col sep=comma, x=s, y=l_max, discard if not={mode}{miku}] {\datapath sl.csv};
\addplot[fill=lightgreen!35, draw=none, forget plot] fill between[of=lmin_miku and lmax_miku];
% 路径（blocked 时按 blocked_s 截断）
\addplot[vibrantgreen, line width=2pt, unbounded coords=discard] table[col sep=comma, x=s, y=l_path, discard if not={mode}{miku}] {\datapath sl.csv};
\addlegendentry{MIKU}


% 障碍物 水马L1
\fill[dangerred!70, draw=dangerred, thick] (axis cs:19.600,1.550) rectangle (axis cs:20.400,1.950);
\node[anchor=south, font=\tiny, color=dangerred!80!black] at (axis cs:20.000,2.250) {水马L1};
% 障碍物 水马L2
\fill[dangerred!70, draw=dangerred, thick] (axis cs:29.600,1.550) rectangle (axis cs:30.400,1.950);
\node[anchor=south, font=\tiny, color=dangerred!80!black] at (axis cs:30.000,2.250) {水马L2};
% 障碍物 水马L3
\fill[dangerred!70, draw=dangerred, thick] (axis cs:39.600,1.550) rectangle (axis cs:40.400,1.950);
\node[anchor=south, font=\tiny, color=dangerred!80!black] at (axis cs:40.000,2.250) {水马L3};
% 障碍物 锥R1
\fill[dangerred!70, draw=dangerred, thick] (axis cs:19.750,-1.425) rectangle (axis cs:20.250,-1.275);
\node[anchor=south, font=\tiny, color=dangerred!80!black] at (axis cs:20.000,-0.975) {锥R1};
% 障碍物 锥R2
\fill[dangerred!70, draw=dangerred, thick] (axis cs:29.750,-1.425) rectangle (axis cs:30.250,-1.275);
\node[anchor=south, font=\tiny, color=dangerred!80!black] at (axis cs:30.000,-0.975) {锥R2};
% 障碍物 锥R3
\fill[dangerred!70, draw=dangerred, thick] (axis cs:39.750,-1.425) rectangle (axis cs:40.250,-1.275);
\node[anchor=south, font=\tiny, color=dangerred!80!black] at (axis cs:40.000,-0.975) {锥R3};
\node[rectangle, fill=deepblue!75, draw=deepblue, thick, minimum width=18pt, minimum height=8pt, inner sep=0pt] (egobox) at (axis cs:0.000,0.000) {};
\fill[white] ([xshift=0pt]egobox.east) -- ([xshift=-3pt,yshift=2.5pt]egobox.east) -- ([xshift=-3pt,yshift=-2.5pt]egobox.east) -- cycle;
\node[anchor=south west, font=\tiny\bfseries, color=deepblue] at (axis cs:0.20,1.10) {EGO};
\end{axis}
\end{tikzpicture}
